\documentclass[10pt,a4paper]{amsart}
\usepackage[margin=19mm]{geometry}
\usepackage{amsmath,amssymb,mathtools}
\usepackage[scr=rsfs,cal=euler]{mathalfa}
\usepackage{xfrac}
\usepackage[unicode,hidelinks,bookmarks=false]{hyperref}
\newcommand{\ie}{i.e.\ }
\newcommand{\cf}{cf.\ }
\newcommand{\resp}{respectively\ }
\newcommand{\Subsection}{\subsection}
\providecommand{\nobreakdash}{\nobreak\hbox{-}}
\setlength{\emergencystretch}{2em}
\multlinegap0pt
\usepackage{tikz}
\usetikzlibrary{cd}
\usepackage{amssymb}
\usepackage{tikz}
\usepackage{xcolor}
\usepackage{dsfont}
\usepackage{mathtools}
\usepackage{enumerate}
\usepackage[shortlabels]{enumitem}
\newtheorem{proposition}{Proposition}
\title{Torsion primes for spaces of commuting elements in Lie groups}
\author{Simon Gritschacher, Chris Hone}
\date{Source: arXiv:2608.21191v1 (CC BY 4.0)}
\begin{document}
\maketitle

\noindent\emph{Source and licence.} Torsion primes for spaces of commuting elements in Lie groups. Version arXiv:2608.21191v1 (CC BY 4.0), distributed by arXiv under the Creative Commons Attribution 4.0 International licence, http://creativecommons.org/licenses/by/4.0/, as recorded in the arXiv licence metadata for this exact version and shown on the abstract page https://arxiv.org/abs/2608.21191v1. Full licence notice: \texttt{licenses/arxiv-2608.21191-CC-BY-4.0.txt}.\par\smallskip

\noindent\textit{Editorial scope.} Counted unit: the article's proposition prop:ss (source0.tex lines 632--635). The article's complete original proof of it is delivered with it (source0.tex lines 637--672, one \texttt{proof} environment of 36 lines). No mathematics is written, repaired or re-derived: the statement and the proof are the article's own lines, in the article's own notation, and no retained line is substituted or re-flowed.  No further source-local context is needed: every cross-reference of every retained line resolves inside the two delivered spans.  Nothing was omitted from any retained span, so the package carries no omission record.  No retained line contains a citation, so the wrapper carries no bibliography and no bibliography entry.  The retained lines contain the article's own diagram code, which the wrapper typesets with the standard tikz packages; no image file is delivered and the package declares no asset dependency.  Editorial formatting only, in addition: an AMS \texttt{amsart} preamble that restates the article's own declarations and macros used by the retained lines, namely the environments \texttt{proposition} on separate counters, together with the standard packages those lines need.  The retained environments print in this document as Proposition 1 in order of appearance, whereas the article prints the same items with the numbering of its own full document.  The complete native source of the article as distributed in the arXiv e-print for this exact version is bundled unchanged as \texttt{sources/208229/source0.tex}.
\begin{proposition} \label{prop:ss}
Let $F\to E\to B$ be a fibre bundle over a connected base, whose fibre $F$ is a finite CW-complex of dimension $f$. Assume there exist nested closed sub-bundles \[E_{\leq 0}\subset E_{\leq 1}\subset \cdots \subset E_{\leq f}=E\] such that the fibres $F_{\leq k}$ of $E_{\leq k}$ are finite CW-complexes of dimension $\leq k$. If for some coefficients $R$, and all $0\leq k\leq f$ the induced maps\[H^i(F,R)\to H^i(F_{\leq k},R)\] are isomorphisms in all degrees $i \leq k$, then the Serre spectral sequence degenerates at $E_2$ and the associated filtration splits. In particular, we have an isomorphism\[H^n(E,R)\cong \bigoplus_{p+q=n} H^p(B,\mathcal{H}^q(F,R))\]
for every $n\geq 0$.
\end{proposition}

\begin{proof}
We use $R$-coefficients throughout the proof. First we prove collapse of the spectral sequence for $E_{\leq k}$ for every $k$. We induct over the finite sequence of sub-bundles $E_{\leq k}\subset E$, noting that the base case $E_{\leq 0}$ follows at once as the $E_2$-page lies on a single line.

Assume the statement holds for $E_{\leq k-1}$. The inclusion of bundles \[E_{\leq k-1}\to E_{\leq k}\] induces a morphism of Serre spectral sequences \[\phi:E^{p,q}_{\bullet}\to \widehat{E}^{p,q}_{\bullet}.\] By induction, we know that all the differentials in the target spectral sequence are zero. Our assumption on the fibre cohomology implies that, for $q\leq k-1$, the induced map of $R$ local systems on $B$ \[\mathcal{H}^q(F_{\leq k})\to \mathcal{H}^q(F_{\leq k-1})\] is an isomorphism, and $\mathcal{H}^q(F_{\leq k})=0$ for $q> k$. Consider the first page $E_l$, $l\geq 2$, with a potentially nonzero differential, which is therefore of the form \[d:E_l^{p,k}\to E_l^{p+l,k+1-l}.\]
By naturality of the Serre spectral sequence, we have a commutative diagram induced by $\phi$:
\begin{equation}\label{Diagram:nat of SS}
    \begin{tikzcd}
E^{p,k}_{l}\arrow[d,"d"]\arrow[rr,"\phi_l^{p,k}"]&& \widehat{E}^{p,k}_{l}\arrow[d,"0"]\\
E^{p+l,k+1-l}_{l}\arrow[rr,"\phi_l^{p+l,k+1-l}"]&& \widehat{E}^{p+l,k+1-l}_{l}.
\end{tikzcd}
\end{equation}
By the minimality of $l\geq 2$, we see that $\phi_l^{p+l,k+1-l}$ equals $\phi_2^{p+l,k+1-l}$, which we may identify with the canonical map
 \[H^{p+l}(B,\mathcal{H}^{k+1-l}(F_{\leq k}))\to H^{p+l}(B,\mathcal{H}^{k+1-l}(F_{\leq k-1}))\, .\]
This map is an isomorphism, since $k+1-l\leq k-1$. By commutativity of (\ref{Diagram:nat of SS}), we deduce that $d$ must be zero, and so the spectral sequence collapses at $E_2$.
 
 
Now we prove the second statement of the proposition, namely that the associated filtration splits. For $0\leq k \leq f$, let
\[
0\subset G^nH^n(E_{\leq k})\subset \cdots \subset G^{1}H^n(E_{\leq k})\subset G^{0}H^n(E_{\leq k})=H^n(E_{\leq k})
\]
denote the finite, descending filtration of $H^n(E_{\leq k})$ to which the Serre spectral sequence converges. For every $p\geq 0$ we have
 \[G^pH^n(E_{\leq k})/G^{p+1}H^n(E_{\leq k})\cong E^{p,n-p}_\infty=E_2^{p,n-p}\cong H^p(B,\mathcal{H}^{n-p}(F_{\leq k})).\]
By considering the associated graded, we see that the map \[G^pH^n(E_{\leq k})\to G^pH^n(E_{\leq k-1})\] is an isomorphism if $p\geq n-k+1$.

Now we show that for every $0\leq p\leq n-1$ the inclusion
\[
G^{p+1}H^n(E)\to G^{p}H^n(E)
\]
is split. By naturality of the Serre spectral sequence, the bundle maps $E_{\leq n-p-1}\to E_{\leq n-p}\to E$ (set $E_{\leq k}:=E$ if $k\geq f$) induce a commutative diagram:
 \[\begin{tikzcd}
    G^{p+1}H^n(E)\arrow[d,"\simeq"]\arrow[r]&G^{p}H^n(E)\arrow[d,"\simeq"]\\
     G^{p+1}H^n(E_{\leq n-p})\arrow[r]\arrow[d,"\simeq"]& G^{p}H^n(E_{\leq n-p})\arrow[d]\\
     G^{p+1}H^n(E_{\leq n-p-1})\arrow[r,"\simeq"]& G^{p}H^n(E_{\leq n-p-1})\, .&
 \end{tikzcd}\]
The lower horizontal map is an isomorphism, because the cohomology of the fibre $F_{\leq n-p-1}$ vanishes in degrees $\geq n-p$. The composition provides the desired splitting.
\end{proof}

\end{document}
